\subsection{点分布的张量积}

13.4.1. 设 \(\mathbf{X}_1\) 与 \(\mathbf{X}_2\) 是两个 \(\mathbf{C}^r\) 类流形，设 \(x_1 \in \mathbf{X}_1\)，\(x_2 \in \mathbf{X}_2\)，以及 \(t_1 \in \mathrm{T}_{x_1}^{(k_1)}(\mathbf{X}_1)\)，\(t_2 \in \mathrm{T}_{x_2}^{(k_2)}(\mathbf{X}_2)\)，其中 \(k_1 + k_2 \leqslant r\)。令 \(\mathbf{X} = \mathbf{X}_1 \times \mathbf{X}_2\)，\(x = (x_1, x_2)\)，\(k = k_{1} + k_{2}\)。对 \(i = 1, 2\)，设 \(c_{i} = (U_{i}, \varphi_{i}, E_{i})\) 是 \(X_{i}\) 的以 \(x_{i}\) 为中心的卡，并用 \(\tilde{t}_{i}\) 表示 \(TS(E_{i})\) 中满足 \(\theta_{c_i}(\tilde{t}_{i}) = t_{i}\) 的元素（参见 13.2.1）。设 \(\sigma\) 是 \(TS(E_{1}) \otimes TS(E_{2})\) 到 \(TS(E_{1} \times E_{2})\) 上的典范同构（参见 A, IV, § 5, no. 5）。元素 \(\sigma(\tilde{t}_{1} \otimes \tilde{t}_{2})\) 是 \(\tilde{t}_{1}\) 与 \(\tilde{t}_{2}\)（视为 \(TS(E_{1} \times E_{2})\) 中的元素，借助典范单射 \(TS(E_{i}) \to TS(E_{1} \times E_{2})\)）的对称积；它属于 \(TS^{(k)}(E_{1} \times E_{2})\)。令 \(c = c_{1} \times c_{2}\)；它是 X 的以 \(x\) 为中心的卡。\(\theta_{c}\) 把 \(\sigma(\tilde{t}_{1} \otimes \tilde{t}_{2})\) 映为 \(T_{x}^{(k)}(X)\) 的一个元素，它与卡 \(c_{i}\) 的选取无关。把它称为 \(t_{1}\) 与 \(t_{2}\) 的直积或对称张量积（甚至简称张量积），并记作 \(t_{1} \times t_{2}\) 或 \(t_{1} \otimes t_{2}\)。

我们有 \(\varepsilon_{x_1} \otimes \varepsilon_{x_2} = \varepsilon_x\)。\(t_1 \otimes t_2\) 的常数项是 \(t_1\) 与 \(t_2\) 的常数项之积。

同构 \((y_{1}, y_{2}) \mapsto (y_{2}, y_{1})\)（从 \(\mathbf{X}_{1} \times \mathbf{X}_{2}\) 到 \(\mathbf{X}_{2} \times \mathbf{X}_{1}\) 上的同构）把 \(t_{1} \otimes t_{2}\) 变为 \(t_{2} \otimes t_{1}\)。

13.4.2（“结合性与函子性”）。设 \(X_{i}\)（i = 1, 2, 3）是 \(C^{r}\) 类流形，并设 \(x_{i} \in X_{i}\)，\(t_{i} \in \mathbf{T}_{x_{i}}^{(k_{i})}(\mathbf{X}_{i})\)，其中 \(k_{1} + k_{2} + k_{3} \leqslant r\)。我们有

\[
\left(t _ {1} \otimes t _ {2}\right) \otimes t _ {3} = t _ {1} \otimes \left(t _ {2} \otimes t _ {3}\right);
\]

把这个点分布记作 \(t_{1} \otimes t_{2} \otimes t_{3}\)。以同样方式定义任意有限的张量积。

设 \(\varphi_{1}\colon\mathbf{X}_{1}\to\mathbf{Y}_{1}\) 与 \(\varphi_{2}\colon\mathbf{X}_{2}\to\mathbf{Y}_{2}\) 是 \(\mathbf{C}^{r}\) 类流形态射，并设 \(t_{1}\in\mathbf{T}^{(k_{1})}(\mathbf{X}_{1})\)，\(t_{2}\in\mathbf{T}^{(k_{2})}(\mathbf{X}_{2})\)，其中 \(k_{1}+k_{2}\leqslant r\)。我们有

\[
(\varphi_ {1} \times \varphi_ {2}) _ {*} (t _ {1} \otimes t _ {2}) = \varphi_ {1 *} (t _ {1}) \otimes \varphi_ {2 *} (t _ {2}).
\]

13.4.3. 沿用 13.4.1 的记号，设 \(F_{1}\)、\(F_{2}\) 与 F 是分离的多范数空间，并设 \((u_{1}, u_{2}) \mapsto u_{1}.u_{2}\) 是从 \(F_{1} \times F_{2}\) 到 F 的一个连续双线性映射。设

\[
f _ {i} \colon \mathrm{X} _ {i} \rightarrow \mathrm{F} _ {i} \quad (i = 1, 2)
\]

是 \(\mathbf{C}^r\) 类映射，并用下式定义 \(f_1 \otimes f_2: \mathbf{X} \to \mathbf{F}\)：

\[
(f _ {1} \otimes f _ {2}) (y _ {1}, y _ {2}) = f _ {1} (y _ {1}). f _ {2} (y _ {2}).
\]

映射 \(f_{1} \otimes f_{2}\) 是 \(\mathbf{C}^{r}\) 类的，且我们有

\[
\langle f _ {1} \otimes f _ {2}, t _ {1} \otimes t _ {2} \rangle = \langle f _ {1}, t _ {1} \rangle . \langle f _ {2}, t _ {2} \rangle .
\]

13.4.4. 沿用 13.4.1 中的记号，设 \(f\) 是从 X 到分离的多范数空间 F 的一个 \(\mathbf{C}^{r}\) 类映射。若 \(y_{1} \in \mathbf{X}_{1}\)，用 \(f_{y_{1}}\) 表示映射 \(y_{2} \mapsto f(y_{1}, y_{2})\)（从 \(\mathbf{X}_{2}\) 到 F 的映射），并令 \(g(y_{1}) = \langle f_{y_{1}}, t_{2} \rangle\)。这样定义的函数 \(g: \mathbf{X}_{1} \to \mathbf{F}\) 是 \(\mathbf{C}^{r - k_{1}}\) 类的，且我们有

\[
\langle f, t _ {1} \otimes t _ {2} \rangle = \langle g, t _ {1} \rangle ,
\]

换言之

\[
\langle f, t _ {1} \otimes t _ {2} \rangle = \langle y _ {1} \mapsto \langle y _ {2} \mapsto f (y _ {1}, y _ {2}), t _ {2} \rangle , t _ {1} \rangle .
\]

同样我们有

\[
\langle f, t _ {1} \otimes t _ {2} \rangle = \langle y _ {2} \mapsto \langle y _ {1} \mapsto f (y _ {1}, y _ {2}), t _ {1} \rangle , t _ {2} \rangle .
\]

13.4.5. 沿用 13.4.1 中的记号，设

\[
\alpha_ {k _ {1}, k _ {2}} \colon \mathrm{T} _ {x _ {1}} ^ {(k _ {1})} (\mathrm{X} _ {1}) \otimes \mathrm{T} _ {x _ {2}} ^ {(k _ {2})} (\mathrm{X} _ {2}) \rightarrow \mathrm{T} _ {x} ^ {(k)} (\mathrm{X})
\]

为由双线性映射 \((t_{1}, t_{2}) \mapsto t_{1} \otimes t_{2}\) 定义的线性映射；该映射是单射；当 \(k_{1}\) 与 \(k_{2}\) 都是无穷时，它是双射。

若 \(n \leqslant r\)，用 \(\mathrm{T}_{x}^{(n)}(\mathrm{X}_{1}, \mathrm{X}_{2})\) 表示 \(\mathrm{T}_{x_{1}}^{(n)}(\mathrm{X}_{1}) \otimes \mathrm{T}_{x_{2}}^{(n)}(\mathrm{X}_{2})\) 中由诸子空间 \(\mathrm{T}_{x_{1}}^{(k_{1})}(\mathrm{X}_{1}) \otimes \mathrm{T}_{x_{2}}^{(k_{2})}(\mathrm{X}_{2})\)（\(k_{1} + k_{2} \leqslant n\)）生成的向量子空间。存在一个线性映射

\[
\alpha_ {n} \colon \mathrm{T} _ {x} ^ {(n)} (\mathrm{X} _ {1}, \mathrm{X} _ {2}) \rightarrow \mathrm{T} _ {x} ^ {(n)} (\mathrm{X})
\]

且是唯一的、延拓上述映射 \(\alpha_{k_1,k_2}\) 的线性映射；它是一个同构。在下文中，我们把 \(\mathrm{T}_x^{(n)}(\mathrm{X})\) 借助 \(\alpha_n^{-1}\) 与 \(\mathrm{T}_x^{(n)}(\mathrm{X}_1,\mathrm{X}_2)\)（\(\mathrm{T}_{x_1}^{(n)}(\mathrm{X}_1)\otimes \mathrm{T}_{x_2}^{(n)}(\mathrm{X}_2)\) 的向量子空间）恒同。

13.4.6. 我们保留 13.4.5 和 13.3.1 的记号。通过取商，\(\alpha_{k_1,k_2}\) 定义了一个线性映射

\[
\varepsilon_ {k _ {1}, k _ {2}} \colon \operatorname{gr} _ {k _ {1}} \mathrm{T} _ {x _ {1}} ^ {(r)} (\mathrm{X} _ {1}) \otimes \operatorname{gr} _ {k _ {2}} \mathrm{T} _ {x _ {2}} ^ {(r)} (\mathrm{X} _ {2}) \rightarrow \operatorname{gr} _ {k} \mathrm{T} _ {x} ^ {(r)} (\mathrm{X}).
\]

诸映射 \(\varepsilon_{k_1,k_2}\)（\(k_1 + k_2 \leqslant r\)）是一个零次分次线性映射的分量

\[
\varepsilon : \bigoplus_ {k _ {1} + k _ {2} \leqslant r} \left(\operatorname{gr} _ {k _ {1}} \mathrm{T} _ {x _ {1}} ^ {(r)} (\mathrm{X} _ {1}) \otimes \operatorname{gr} _ {k _ {2}} \mathrm{T} _ {x _ {2}} ^ {(r)} (\mathrm{X} _ {2})\right)\rightarrow \bigoplus_ {k \leqslant r} \operatorname{gr} _ {k} \mathrm{T} _ {x} ^ {(r)} (\mathrm{X})
\]

且该映射是一个同构。图表

\[
\begin{array}{l} \bigoplus_ {k _ {1} + k _ {2} \leqslant r} (\operatorname{gr} _ {k _ {1}} \mathrm{T} _ {x _ {1}} ^ {(r)} (\mathrm{X} _ {1}) \otimes \operatorname{gr} _ {k _ {2}} \mathrm{T} _ {x _ {2}} ^ {(r)} (\mathrm{X} _ {2})) \xrightarrow {\varepsilon} \operatorname{gr} \mathrm{T} _ {x} ^ {(r)} (\mathrm{X}) \\ \Biggl \downarrow i _ {1 2} \\ \bigoplus_ {k _ {1} + k _ {2} \leqslant r} (\mathrm{TS} ^ {k _ {1}} (\mathrm{T} _ {x _ {1}} (\mathrm{X} _ {1})) \otimes \mathrm{TS} ^ {k _ {2}} (\mathrm{T} _ {x _ {2}} (\mathrm{X} _ {2}))) \xrightarrow {\sigma} \bigoplus_ {k \leqslant r} \mathrm{TS} ^ {k} (\mathrm{T} _ {x _ {1}} (\mathrm{X} _ {1}) \times \mathrm{T} _ {x _ {2}} (\mathrm{X} _ {2})) \end{array}
\]

是交换的（在这个图中，\(i_{12}\) 表示由 \(i_{\mathbf{x}_1} \otimes i_{\mathbf{x}_2}\) 诱导的同态，而 \(\sigma\) 是 A, IV, § 5, no. 5 中定义的同构）。若 \(r\) 是无穷，这个图可以简单地写成：

\[
\begin{array}{l} \operatorname{gr} \mathrm{T} _ {x _ {1}} ^ {(\infty)} (\mathrm{X} _ {1}) \otimes \operatorname{gr} \mathrm{T} _ {x _ {2}} ^ {(\infty)} (\mathrm{X} _ {2}) \xrightarrow {\varepsilon} \operatorname{gr} \mathrm{T} _ {x} ^ {(\infty)} (\mathrm{X}) \\ \Biggl \downarrow_ {i _ {\mathbf {X} _ {1}} \otimes i _ {\mathbf {X} _ {2}}} \\ \operatorname{TS} (\mathrm{T} _ {x _ {1}} (\mathrm{X} _ {1})) \otimes \operatorname{TS} (\mathrm{T} _ {x _ {2}} (\mathrm{X} _ {2})) \xrightarrow {\sigma} \operatorname{TS} (\mathrm{T} _ {x} (\mathrm{X})). \end{array}
\]

